% GENERATED FILE. Edit canonical lessons, then run npm run generate:books.
% canonical-source-hash: fnv1a64:3ab0f4a0a42bd0a0
% canonical-lessons: LA-C538-otiosus, LA-C538-vaco, LA-C538-feriatus, LA-C538-relaxo, LA-C538-remissio

\chapter{Free time}
\label{ch:la-free-time}
\hlchaptermodality{538}

\begin{chapteropening}
\textbf{By the end of this chapter:} \emph{I can talk about leisure and holidays.}

Complete the last lesson of chapter 538: I can talk about leisure and holidays.
\end{chapteropening}

\section[ōtiōsus, ōtiōsa, ōtiōsum]{ōtiōsus, ōtiōsa, ōtiōsum — at leisure, idle}

\label{lesson:LA-C538-otiosus}



\begin{quote}
Before the new one: say the Latin for faithless, then the Latin for to abandon.
\end{quote}

\subsection*{You'll want to know: ōtiōsus, ōtiōsa, ōtiōsum}
\textbf{ōtiōsus, ōtiōsa, ōtiōsum} — \textquotedblleft{}at leisure, idle\textquotedblright{}.

\begin{quote}
Diē fēstō cīvēs ōtiōsī in forō ambulant.
\end{quote}

\begin{quote}
\emph{On the holiday the citizens stroll at leisure in the forum.}
\end{quote}

From \emph{ōtium}: free of business, with no blame meant. English \emph{otiose}.

\begin{cousinweb}
\textbf{ōtiōsus} is built on \emph{ōtium} \textquotedblleft{}leisure\textquotedblright{} with \emph{-ōsus} (EtymDB; Lewis \& Short). English \emph{otiose} comes from it.
\end{cousinweb}

\begin{grammarlens}[title={What you've built}]
One more word for telling, explaining and describing at B1.
\end{grammarlens}

\subsection*{Your turn}
\begin{itemize}
\raggedright
  \item \emph{Say it:} \emph{ōtiōsus, ōtiōsa, ōtiōsum}
  \item \emph{Say it:} \emph{ōtiōsus, ōtiōsa, ōtiōsum}, once more
  \item \emph{Say it:} a sentence of your own with \emph{ōtiōsus, ōtiōsa, ōtiōsum}
  \item \emph{Recall:} say the Latin for faithless, then the Latin for to abandon, then say \emph{ōtiōsus, ōtiōsa, ōtiōsum} again
\end{itemize}

\begin{tcolorbox}[breakable,colback=teal!4,colframe=teal!35!black,title={Before you move on}]
What does ōtiōsus, ōtiōsa, ōtiōsum mean? (\textquotedblleft{}At leisure, idle\textquotedblright{}.) Say it once more.
\end{tcolorbox}
\section[vacō, vacāre]{vacō, vacāre — to be free, have time (for)}

\label{lesson:LA-C538-vaco}



\begin{quote}
Before the new one: say the Latin for to abandon, then the Latin for at leisure.
\end{quote}

\subsection*{You'll want to know: vacō, vacāre}
\textbf{vacō, vacāre} — \textquotedblleft{}to be free, have time (for)\textquotedblright{}.

\begin{quote}
Hodiē vacō; tēcum ambulāre possum.
\end{quote}

\begin{quote}
\emph{Today I am free; I can go for a walk with you.}
\end{quote}

\emph{vacō, vacāre, vacāvī, vacātum}. To be empty, then to be free for something: \emph{vacāre litterīs} — to have time for literature. English \emph{vacant}, \emph{vacation}.

\begin{cousinweb}
\textbf{vacō} is inherited, through Proto-Italic, from a Proto-Indo-European root meaning \textquotedblleft{}empty\textquotedblright{} (EtymDB). English \emph{vacant} and \emph{vacation} come from it.
\end{cousinweb}

\begin{grammarlens}[title={What you've built}]
One more word for telling, explaining and describing at B1.
\end{grammarlens}

\subsection*{Your turn}
\begin{itemize}
\raggedright
  \item \emph{Say it:} \emph{vacō, vacāre}
  \item \emph{Say it:} \emph{vacō, vacāre}, once more
  \item \emph{Say it:} a sentence of your own with \emph{vacō, vacāre}
  \item \emph{Recall:} say the Latin for to abandon, then the Latin for at leisure, then say \emph{vacō, vacāre} again
\end{itemize}

\begin{tcolorbox}[breakable,colback=teal!4,colframe=teal!35!black,title={Before you move on}]
What does vacō, vacāre mean? (\textquotedblleft{}To be free, have time (for)\textquotedblright{}.) Say it once more.
\end{tcolorbox}
\section[fēriātus, fēriāta, fēriātum]{fēriātus, fēriāta, fēriātum — on holiday, keeping holiday}

\label{lesson:LA-C538-feriatus}



\begin{quote}
Before the new one: say the Latin for at leisure, then the Latin for to be free.
\end{quote}

\subsection*{You'll want to know: fēriātus, fēriāta, fēriātum}
\textbf{fēriātus, fēriāta, fēriātum} — \textquotedblleft{}on holiday, keeping holiday\textquotedblright{}.

\begin{quote}
Puerī fēriātī tōtum diem lūdunt.
\end{quote}

\begin{quote}
\emph{The boys on holiday play all day.}
\end{quote}

The participle of \emph{fērior} — to keep holiday, from \emph{fēriae}, already known.

\begin{cousinweb}
\textbf{fēriātus} is the participle of \emph{fērior} \textquotedblleft{}to keep holiday\textquotedblright{}, built on \emph{fēriae} \textquotedblleft{}holidays\textquotedblright{} (Lewis \& Short).
\end{cousinweb}

\begin{grammarlens}[title={What you've built}]
One more word for telling, explaining and describing at B1.
\end{grammarlens}

\subsection*{Your turn}
\begin{itemize}
\raggedright
  \item \emph{Say it:} \emph{fēriātus, fēriāta, fēriātum}
  \item \emph{Say it:} \emph{fēriātus, fēriāta, fēriātum}, once more
  \item \emph{Say it:} a sentence of your own with \emph{fēriātus, fēriāta, fēriātum}
  \item \emph{Recall:} say the Latin for at leisure, then the Latin for to be free, then say \emph{fēriātus, fēriāta, fēriātum} again
\end{itemize}

\begin{tcolorbox}[breakable,colback=teal!4,colframe=teal!35!black,title={Before you move on}]
What does fēriātus, fēriāta, fēriātum mean? (\textquotedblleft{}On holiday, keeping holiday\textquotedblright{}.) Say it once more.
\end{tcolorbox}
\section[relaxō, relaxāre]{relaxō, relaxāre — to relax, ease; to loosen}

\label{lesson:LA-C538-relaxo}



\begin{quote}
Before the new one: say the Latin for to be free, then the Latin for on holiday.
\end{quote}

\subsection*{You'll want to know: relaxō, relaxāre}
\textbf{relaxō, relaxāre} — \textquotedblleft{}to relax, ease; to loosen\textquotedblright{}.

\begin{quote}
Post labōrem animum mūsicā relaxō.
\end{quote}

\begin{quote}
\emph{After work I relax my mind with music.}
\end{quote}

\emph{relaxō, relaxāre, relaxāvī, relaxātum}: \emph{re-} + \emph{laxō} 'to loosen'. English \emph{relax}.

\begin{cousinweb}
\textbf{relaxō} is built on \emph{re-} and \emph{laxō} \textquotedblleft{}to loosen\textquotedblright{}, from \emph{laxus} \textquotedblleft{}slack\textquotedblright{} (EtymDB; Lewis \& Short). English \emph{relax} and \emph{release} come from it.
\end{cousinweb}

\begin{grammarlens}[title={What you've built}]
One more word for telling, explaining and describing at B1.
\end{grammarlens}

\subsection*{Your turn}
\begin{itemize}
\raggedright
  \item \emph{Say it:} \emph{relaxō, relaxāre}
  \item \emph{Say it:} \emph{relaxō, relaxāre}, once more
  \item \emph{Say it:} a sentence of your own with \emph{relaxō, relaxāre}
  \item \emph{Recall:} say the Latin for to be free, then the Latin for on holiday, then say \emph{relaxō, relaxāre} again
\end{itemize}

\begin{tcolorbox}[breakable,colback=teal!4,colframe=teal!35!black,title={Before you move on}]
What does relaxō, relaxāre mean? (\textquotedblleft{}To relax, ease; to loosen\textquotedblright{}.) Say it once more.
\end{tcolorbox}
\section[remissiō, remissiōnis]{remissiō, remissiōnis — relaxation; a letting go}

\label{lesson:LA-C538-remissio}



\begin{quote}
Before the new one: say the Latin for on holiday, then the Latin for to relax.
\end{quote}

\subsection*{You'll want to know: remissiō, remissiōnis}
\textbf{remissiō, remissiōnis} — \textquotedblleft{}relaxation; a letting go\textquotedblright{}.

\begin{quote}
Post studia remissiō animī necessāria est.
\end{quote}

\begin{quote}
\emph{After study, relaxation of the mind is necessary.}
\end{quote}

Feminine, third declension, from \emph{remittō} — to send back, slacken. English \emph{remission}.

\begin{cousinweb}
\textbf{remissiō} is the action noun of \emph{remittō} \textquotedblleft{}to send back, slacken\textquotedblright{}, built on \emph{re-} and \emph{mittō} (Lewis \& Short). English \emph{remission} comes from it.
\end{cousinweb}

\begin{grammarlens}[title={What you've built}]
One more word for telling, explaining and describing at B1.
\end{grammarlens}

\subsection*{Your turn}
\begin{itemize}
\raggedright
  \item \emph{Say it:} \emph{remissiō, remissiōnis}
  \item \emph{Say it:} \emph{remissiō, remissiōnis}, once more
  \item \emph{Say it:} a sentence of your own with \emph{remissiō, remissiōnis}
  \item \emph{Recall:} say the Latin for on holiday, then the Latin for to relax, then say \emph{remissiō, remissiōnis} again
\end{itemize}

\begin{tcolorbox}[breakable,colback=teal!4,colframe=teal!35!black,title={Before you move on}]
What does remissiō, remissiōnis mean? (\textquotedblleft{}Relaxation; a letting go\textquotedblright{}.) Say it once more.
\end{tcolorbox}
